% Generated draft: requires mathematical and source review.
\documentclass[11pt]{article}
\usepackage[margin=25mm]{geometry}
\usepackage{fontspec}
\setmainfont{Latin Modern Roman}
\usepackage{amsmath,amssymb,mathtools,unicode-math}
\setmathfont{Latin Modern Math}
\usepackage{enumitem,needspace,xurl,hyperref}
\setlength{\parindent}{0pt}
\setlength{\parskip}{5pt}
\setlength{\emergencystretch}{3em}
\newcommand{\N}{\mathbb{N}}
\newcommand{\Z}{\mathbb{Z}}
\newcommand{\Q}{\mathbb{Q}}
\newcommand{\R}{\mathbb{R}}
\newcommand{\C}{\mathbb{C}}
\newcommand{\F}{\mathbb{F}}
\newcommand{\A}{\mathbb{A}}
\newcommand{\PP}{\mathbb{P}}
\newcommand{\E}{\mathbb{E}}
\newcommand{\eps}{\varepsilon}
\newcommand{\dd}{\,\mathrm{d}}
\begin{document}
\section*{UnsolvedMath ID: 30006169}
\textbf{Source problem code:} catalogue entry
\subsection{Definitions and mathematical statement}
Work in a ground universe \(V\). Let \(\omega<\kappa<\lambda<\mu\) be regular cardinals of \(V\), and assume the source hypothesis \(\operatorname{LIP}(\lambda,\mu)\). The supplied record does not define \(\operatorname{LIP}\); it is retained as an unresolved hypothesis, without assigning it a meaning.

For a forcing notion \(\mathbb P\in V\), say that \(\mathbb P\) preserves every cardinal at most \(\lambda\) if
\[
\forall\nu\,\bigl((\nu\leq\lambda\ \text{and}\ V\models\text{``}\nu\text{ is a cardinal''})\ \Longrightarrow\ \mathbb P\Vdash\text{``}\check\nu\text{ is a cardinal''}\bigr).
\]
Here \(\check\nu\) is the canonical forcing name for the ground-model ordinal \(\nu\). Equivalently, every such ground-model cardinal remains a cardinal in every \(\mathbb P\)-generic extension.

\emph{Question, conditional on the missing source definitions.} Under the hypotheses above, does there exist a Namba-type minimal forcing notion \(\mathbb P\) that preserves every cardinal at most \(\lambda\)? The original formulation asks for a forcing ``like the one used in Theorem 1.'' That theorem and its forcing construction are absent from the supplied record. Consequently, the intended class of Namba-type forcings, the meaning of minimality, the required change of cofinality or collapse behavior, and the role of \(\kappa\) cannot be specified from this record. This is a schematic statement requiring those definitions before it becomes a fully determined mathematical question.

\emph{Source qualifications and surrounding discussion.} The source says an affirmative answer seems plausible, but supplies no proof. It also discusses questions about cardinal collapses and embeddings of forcing orders for Namba-style forcings, including questions attributed to Bukovsk\'y and Copl\'akov\'a and questions at the end of its reference [6]. It emphasizes that such questions depend on the dimensions and splitting behavior of the forcing, and asks how far the main theorem can be stratified. Neither that main theorem nor reference [6] is identified in the supplied excerpt.

The excerpt additionally considers the conditional situation \(\mu=\nu^{++}\), with \(\mu\) not measurable in any inner model, and asserts that there would then be a sequence of length \(\mu^+\) witnessing the non-saturation of \(\mu\cap\operatorname{cof}(\nu^+)\). Here \(\nu^{++}\) denotes the second successor cardinal of \(\nu\), and \(\mu\cap\operatorname{cof}(\nu^+)=\{\alpha<\mu:\operatorname{cf}(\alpha)=\nu^+\}\), read in the ground universe. The relevant ideal and the precise witnessing property are not specified. The excerpt suggests using ideas from a proof attributed to Shelah concerning singularizing \(\mu\) while collapsing \(\mu^+\), but does not supply its hypotheses or citation and explicitly says that the role of minimality in such an argument is unclear. These remarks are retained as source discussion, not as additional hypotheses or verified assertions; the auxiliary \(\nu\) in this discussion is distinct from the universally quantified cardinal in the preservation requirement.
\subsection{Short English statement}
The supplied question asks whether, assuming \(\operatorname{LIP}(\lambda,\mu)\) for regular cardinals \(\omega<\kappa<\lambda<\mu\), a suitable minimal Namba-type forcing can preserve every ground-model cardinal at most \(\lambda\). The hypothesis, the intended forcing construction, and minimality are not defined in the record, so the statement remains conditional on recovering them. The record labels the problem open; this draft does not independently verify that status.
\subsection{Sources}
\begin{itemize}
\item Local UnsolvedMath registry: \texttt{data/unsolvedmath/problems.json}, with \texttt{catalog/problems.json} as the supplementary registry; permanent UnsolvedMath ID \(30006169\), title \emph{Cardinal-Preserving Namba-Type Forcing}, exact source code \texttt{catalogue entry}. This draft uses the supplied complete record; the registry files were not independently inspected.
\item Supplied report provenance: \emph{Set Theory} (2025), Oberwolfach Reports, pp.~75--140, DOI \texttt{10.4171/owr/2025/2}. \url{https://doi.org/10.4171/owr/2025/2}. The record gives the proposed year as 2025 and names no individual proposer.
\item Supplied literature source: Levine, \emph{On Namba Forcing and Minimal Collapses}. \url{https://resolve.cambridge.org/core/services/aop-cambridge-core/content/view/139989528588848BDDFD80F18AF63ED9/S2050509425101060a.pdf/on-namba-forcing-and-minimal-collapses.pdf}.
\item Supplied literature source: \emph{Oberwolfach Report 2/2025}. \url{https://content.ems.press/assets/public/full-texts/serials/owr/22/1/14299082/online/10.4171-owr-2025-2.pdf}.
\item Supplied literature assessment, checked on exactly \texttt{2026-08-22}: the record labels the problem open and states that Levine's paper retains the cardinal-preservation question as an explicit open question, without a general preservation theorem resolving it. This is supplied dated curation, described as literature triage derived from the supplied Oberwolfach curation; neither its status assessment nor its cited claims have been independently verified here.
\item The registry marks the statement \texttt{corrected\_verified} and reports that it was reconstructed and checked against the official Oberwolfach report. This is supplied verification metadata, not independent verification by this draft. The excerpt's references to Theorem 1, the main theorem, reference [6], Bukovsk\'y and Copl\'akov\'a, and a proof of Shelah lack complete identifying information in the supplied record.
\end{itemize}
\end{document}
